\documentclass[letterpaper,10pt,journal,twoside]{IEEEtran}

% =========================
% Packages
% =========================
\usepackage{amsmath,amsfonts,amssymb}
\usepackage{bm}
\usepackage[mathscr]{euscript}

\usepackage{algorithm}
\usepackage{algorithmic}
\usepackage{array}
\usepackage{booktabs}
\usepackage{multirow}
\usepackage{arydshln}

\usepackage{graphicx}
\usepackage[dvipsnames]{xcolor}
\usepackage[caption=false,font=normalsize,labelfont=sf,textfont=sf]{subfig}
\usepackage[skip=1.5pt]{caption}

\usepackage{textcomp}
\usepackage{stfloats}
\usepackage{url}
\usepackage{verbatim}
\usepackage{cite}

% Chinese support.
% Use fandol for portability across machines. If your system has Windows CJK
% fonts correctly installed, this can be changed to fontset=windows.
\usepackage[UTF8,scheme=plain,fontset=fandol]{ctex}

% Times-like text and math fonts, commonly used with IEEEtran.
\usepackage{newtxtext,newtxmath}

% =========================
% Layout tuning
% =========================
\setlength{\floatsep}{5pt}
\setlength{\textfloatsep}{5pt}
\setlength{\intextsep}{5pt}

\hyphenation{op-tical net-works semi-conduc-tor IEEE-Xplore}

% =========================
% Theorem environments
% =========================
\newtheorem{definition}{Definition}
\newtheorem{theorem}{Theorem}
\newtheorem{remark}{Remark}
\newtheorem{corollary}{Corollary}
\newtheorem{proof}{Proof}

% =========================
% Custom commands
% =========================
\definecolor{r}{rgb}{1,0,0}
\definecolor{b}{rgb}{0,0,1}

\newcommand{\todo}[1]{\textcolor{r}{[TODO: #1]}}
\newcommand{\revise}[1]{\textcolor{b}{#1}}

\begin{document}

% =========================
% Title and authors
% =========================
\title{Paper Title}

\author{Author~A,
        Author~B,
        and~Author~C%
\thanks{This work was supported by xxx. Corresponding author: xxx.}
\thanks{Author A, Author B, and Author C are with xxx.
Email: author@example.com.}
}

\maketitle

% =========================
% Abstract and keywords
% =========================
\begin{abstract}
Write a concise summary of the problem, method, main contribution, and results.
For IEEE-style papers, the abstract is usually one paragraph and should avoid
citations, equations, and undefined abbreviations when possible.
\end{abstract}

\begin{IEEEkeywords}
Keyword one, keyword two, keyword three, keyword four.
\end{IEEEkeywords}

% =========================
% Main text
% =========================
\section{Introduction}
\label{sec:introduction}

Introduce the background, motivation, research gap, objective, main idea, and
contributions of the paper. A common structure is: why the topic matters, what
has been done, what remains unsolved, what this paper does, and why the results
are useful.

The main contributions of this paper are summarized as follows:
\begin{itemize}
    \item We identify or formulate an important problem in the target domain.
    \item We propose a method, framework, analysis, system, or design to address
    the problem.
    \item We validate the proposed approach through experiments, case studies,
    theoretical analysis, or comparative evaluation.
\end{itemize}

The remainder of this paper is organized as follows.
Section~\ref{sec:background} introduces the necessary background and problem
definition. Section~\ref{sec:related} reviews related work.
Section~\ref{sec:approach} presents the proposed approach.
Section~\ref{sec:evaluation} reports the evaluation results.
Section~\ref{sec:conclusion} concludes this paper.

\section{Background and Problem Definition}
\label{sec:background}

Introduce the necessary concepts, notation, assumptions, and task definition.
This section can also be named ``Preliminaries'', ``System Model'',
``Problem Formulation'', or ``Materials and Methods'', depending on the paper
type.

\begin{equation}
    y = f(x;\theta).
    \label{eq:example}
\end{equation}

Use equations such as Eq.~\eqref{eq:example} to define important quantities,
models, objectives, or constraints.

\section{Related Work}
\label{sec:related}

Review representative work and explain the difference between existing methods
and your work. This section may be organized by research direction, technical
category, application scenario, or historical development. Add citations with
\verb|\cite{}| after adding entries to \verb|ref.bib|.

\section{Proposed Approach}
\label{sec:approach}

Describe the central contribution of the paper. Depending on the topic, this
section can present a method, theory, model, algorithm, system, architecture,
workflow, design principle, or analytical framework.

\subsection{Overview}
\label{subsec:framework}

Give a high-level overview before entering details. A figure is often useful
here to show the relationship between inputs, modules, outputs, and evaluation.

\subsection{Detailed Design}
\label{subsec:design}

Describe the proposed components, assumptions, derivations, implementation
details, or design choices. Split this subsection into multiple subsections if
the paper has several major parts.

\subsection{Procedure}
\label{subsec:procedure}

\begin{algorithm}[t]
\caption{General Procedure}
\label{alg:procedure}
\begin{algorithmic}[1]
\STATE Prepare the input data, materials, assumptions, or system state.
\STATE Initialize the required variables, modules, or parameters.
\FOR{each iteration, case, sample, or processing step}
    \STATE Apply the proposed operation, analysis, or decision rule.
    \STATE Record intermediate results if needed.
\ENDFOR
\STATE Return the final output, result, or conclusion.
\end{algorithmic}
\end{algorithm}

\section{Evaluation}
\label{sec:evaluation}

Report how the proposed work is validated. This can include experiments,
simulations, case studies, benchmarks, user studies, theoretical proofs, error
analysis, or practical deployment results.

\subsection{Setup}
\label{subsec:setup}

Describe the data, materials, apparatus, environment, implementation details,
baselines, comparison methods, evaluation metrics, or analysis protocol.

\subsection{Results}
\label{subsec:results}

\begin{table}[t]
\centering
\caption{Summary of representative results.}
\label{tab:main_results}
\begin{tabular}{lccc}
\toprule
Item & Metric 1 & Metric 2 & Metric 3 \\
\midrule
Case A & -- & -- & -- \\
Case B & -- & -- & -- \\
Case C & -- & -- & -- \\
\bottomrule
\end{tabular}
\end{table}

\begin{figure}[t]
\centering
% \includegraphics[width=\linewidth]{image/example.pdf}
\caption{Example figure caption.}
\label{fig:example}
\end{figure}

\subsection{Analysis}
\label{subsec:analysis}

Analyze the meaning of the results, compare against prior work, explain
observed trends, and discuss whether the evidence supports the claims.

\subsection{Discussion}
\label{subsec:discussion}

Discuss limitations, threats to validity, practical implications, generality,
failure cases, or future extensions.

\section{Conclusion}
\label{sec:conclusion}

Summarize the paper and briefly mention future work.

% =========================
% References
% =========================
% Add references to ref.bib and cite them in the text with \cite{}.
% If you want to print all entries in ref.bib while drafting, uncomment:
% \nocite{*}
\bibliographystyle{IEEEtran}
\bibliography{ref}

\end{document}
